\chapter*{CONCLUSÃO} 
\addcontentsline{toc}{chapter}{CONCLUSÃO}
\thispagestyle{plain}

Nesta seção são descritas claramente as conclusões retiradas das
discussões e dos experimentos realizados no decorrer da pesquisa, e finalizada a parte textual do trabalho. A Conclusão deve dialogar com a Introdução, melhor dizendo, apontar ou sugerir questões e problemas decorrentes do assunto apresentado, mostrando se os objetivos foram atingidos ou não e se as hipóteses foram ou não confirmadas. Também deve trazer recomendações e uma breve explanação sobre trabalhos futuros.

%MAIS DICAS
%
%Sua conclusão pode apresentar uma estrutura como a seguinte:
%Este trabalho apresentou .... De uma maneira geral a 
%
%A primeira etapa foi ..
%
%
%Para a segunda etapa, que é a ... O algoritmo que separa em regiões foi adaptado ao processamento morfológico resultando em .. 
%
%A terceira e última etapa, que é a ... 
%
%
%
%
%Portanto, as principais contribuições do presente trabalho são: ...
%
%
%E como trabalhos futuros, será implementado um novo método baseado em \textit{roughness index 3D}; serão desenvolvidas redes convolucionais para ..